\hypertarget{appendix-0b3be18411f9}{}\label{referencestart:appendix-0b3be18411f9}
\section{Eigenvalues and eigenspaces}\label{appendix-0b3be18411f9-eigenvalues-and-eigenspaces}

An eigenvector is a nonzero direction that a square linear map scales by its eigenvalue. The eigenspace is the null space of \(A-\lambda I\). Real symmetric matrices have an orthonormal eigenbasis; general matrices may be defective or have sensitive eigenvectors.

\textbf{Use it when:} studying repeated powers, invariant directions, covariance, stability, or matrix functions.

\textbf{Developed in Phase 8:} Eigenanalysis begins with invariant directions, then distinguishes algebraic multiplicity from eigenspace dimension and shows why a defective matrix may not have a full eigenbasis. Power and inverse iteration target selected eigenpairs; QR iteration approaches a Schur form, while residual size and eigenvector sensitivity answer different accuracy questions. \href{../02-Mathematics-for-Visual-Computing/main.pdf\#nameddest=01-foundations-i-ch020-eigenvalues-eigenspaces-and-the-two-kinds-of-multiplicity}{eigenspaces and multiplicity (Mathematics for Visual Computing, Mathematics for Visual Computing, Chapter 4)} · \href{../02-Mathematics-for-Visual-Computing/main.pdf\#nameddest=01-foundations-i-ch020-qr-iteration-and-jacobi-rotations-preserve-the-spectrum-while-simplifying-shape}{eigenvalue iteration and Schur form (Mathematics for Visual Computing, Mathematics for Visual Computing, Chapter 4)}

\subsection{Prerequisites}\label{appendix-0b3be18411f9-prerequisites}

\href{../02-Mathematics-for-Visual-Computing/main.pdf\#nameddest=appendix-bd39c3bf5d64}{Matrices as linear maps and basis images (Mathematics for Visual Computing, Reference notes: Matrices as linear maps and basis images)}, \hyperref[appendix-8572a4a060c1]{Linear systems and fundamental subspaces}, \hyperref[appendix-175ca5ff874d]{Matrix transpose symmetry and trace}. - \hyperref[appendix-ab31ef447dbf]{Matrix multiplication and composition order}

\subsection{Used by}\label{appendix-0b3be18411f9-used-by}

\href{../04-Light-Transport-and-Image-Synthesis/main.pdf\#nameddest=appendix-1fad5e7f36fa}{Singular value decomposition (Light Transport and Image Synthesis, Reference notes: Singular value decomposition)}, \hyperref[appendix-ad175ef3c5bf]{PCA and covariance eigensystems}, \hyperref[appendix-73f5bcf7a462]{Polar decomposition}, \hyperref[appendix-0f4bc2bf2b1f]{Matrix functions}.

\textbf{Teaching path:} \href{../02-Mathematics-for-Visual-Computing/main.pdf\#nameddest=01-foundations-i-ch020-eigenvectors-reveal-the-directions-a-square-map-does-not-turn}{eigenanalysis (Mathematics for Visual Computing, Mathematics for Visual Computing, Chapter 4)}. - \hyperref[appendix-e8ab9836882e]{ODE initial-value problems and stiffness}

\subsection{Full teaching explanation}\label{appendix-0b3be18411f9-full-teaching-explanation}

\begin{itemize}
\tightlist
\item
  \href{../02-Mathematics-for-Visual-Computing/main.pdf\#nameddest=02-foundations-ii-ch025-odes-and-numerical-simulation}{Phase 9 teaching section (Mathematics for Visual Computing, Mathematics for Visual Computing, Chapter 9)}
\item
  \href{../01-Bits-Programs-and-Reliable-Systems/main.pdf\#nameddest=appendix-adf48f5dc5c1}{Visual Recall/09 - Calculus and Optimization - Visual Recall.md (Bits, Programs, and Reliable Systems, Reference notes: 09 - Calculus and Optimization - Visual Recall)} · Phase 9 visual recall
\end{itemize}
